% SemEval-2027 Task 1 (RETECO) system description — CUET_AL_Masaar.
% DRAFT started 2026-09-23. Every number in this file is copied from results/ledger.csv;
% the ledger run_id is given in a comment next to it. \pending{} marks numbers that do not
% exist yet (dev check, test submission). Do not fill one without a ledger row.
\documentclass[11pt]{article}
\usepackage[review]{acl}
\usepackage{times}
\usepackage{latexsym}
\usepackage[T1]{fontenc}
\usepackage[utf8]{inputenc}
\usepackage{microtype}
\usepackage{booktabs}
\usepackage{xcolor}
\newcommand{\pending}[1]{\textcolor{red}{[PENDING: #1]}}

\title{CUET\_AL\_Masaar at SemEval-2027 Task 1: Candidate-Pool Union of Decomposed Steps\\
and Listwise Reasoning Reranking for Temporal Retrieval}

\author{Rafiur Rahman \\
  Chittagong University of Engineering and Technology \\
  \texttt{u2303064@student.cuet.ac.bd}}

\begin{document}
\maketitle

\begin{abstract}
We describe our system for Track~1 of RETECO (temporal retrieval, sub-tasks 1a and 1b). It has
three stages: a 0.6B dense retriever, a fusion of each query's ranking with the rankings of its
decomposed reasoning steps, and a 7B listwise reasoning reranker over the top 30 candidates.
On the training split, macro nDCG@10 over queries goes from 0.0944 (official BM25) to 0.2570
(dense), 0.2685 (step fusion) and, on a random 896-query subset, from 0.2787 to 0.3296 with
reranking. Our main finding concerns step fusion: its gain comes entirely from adding the whole
query's candidates to the step candidates, not from how scores are combined. A parent-query
weight of 0.001 captures the full gain, the choice among sum, max and reciprocal rank fusion
makes no significant difference, and the same union applied in the other direction (adding
the query's candidates to each step) does not help sub-task 1b.
Test results: \pending{official 1a / 1b test nDCG@10 and rank}.
\end{abstract}

\section{Introduction}
RETECO Track~1 is built on TEMPO \citep{abdallah2026tempo}: 13 StackExchange-derived domains,
1,654,055 documents, and questions whose answers depend on temporal reasoning. Sub-task 1a ranks
documents for a whole question; sub-task 1b ranks documents for each of the question's official
decomposed steps. Both are scored by nDCG@10, macro-averaged over queries (1b first averages over
a query's steps). The official BM25 baseline reaches 0.0879 macro over domains on the training
split, which we reproduce exactly on all 13 domains.

We entered 1a and 1b with one pipeline and chose every component on the training split alone,
using domain-stratified bootstrap confidence intervals and paired tests. The development split
was used once, for the frozen system (\S\ref{sec:dev}).

\section{System}
\paragraph{First stage.} \texttt{AQ-MedAI/Diver-Retriever-0.6B} \citep{sun2025diver}, chosen
over its 1.7B and 4B siblings on measured cost: embedding the deduplicated corpus (29.4\% of
documents are byte-identical duplicates) took 6.0 hours on two T4 GPUs in fp16. Queries use the
model's query prompt and documents none; 1b queries use the official template (query, blank
line, \texttt{Step:} instruction). Scores are computed in fp32; computing them in fp16 left
about 70\% of each top-100 in exact-score ties.

\paragraph{Step fusion (1a).} For each question, the whole-query ranking (weight $w$) and each
step ranking (weight 1) are normalised by their maximum score and summed; documents missing from
a list score 0; ties break by corpus order; the top 100 are kept. We use $w{=}0.001$
(\S\ref{sec:fusion}).

\paragraph{Reranking (1a).} ReasonRank-7B \citep{liu2025reasonrank} reranks the fused top 30
with the authors' own prompt and settings: windows of 20 passages moved back to front with step
10, passages truncated to 512 tokens, greedy decoding with a 3,072-token reasoning budget.
Ranks 31--100 are unchanged. We chose it from published BRIGHT results \citep{su2025bright},
where most rerankers below 32B lower the score of a strong retriever and ReasonRank-7B raises it.

\paragraph{1b.} The dense step rankings, without fusion or reranking (\S\ref{sec:fusion}).

\section{Experiments on the training split}
All numbers are macro nDCG@10 over the 1,211 training queries unless stated, with paired
bootstrap 95\% confidence intervals.

\begin{table}[t]
\centering\small
\begin{tabular}{lr}
\toprule
System (1a, train) & nDCG@10 \\
\midrule
Official BM25 & 0.0944 \\ % phase2_bm25_full_r1
Dense (Diver-0.6B) & 0.2570 \\ % phase5 baseline, fp32 scoring
+ step fusion, $w{=}0.001$ & 0.2685 \\ % phase5_stepfuse_w0.001
\midrule
\multicolumn{2}{l}{\emph{random 896-query subset}} \\
+ step fusion & 0.2787 \\ % phase7_rerank_reasonrank7b_d30
+ ReasonRank-7B, top 30 & 0.3296 \\ % phase7_rerank_reasonrank7b_d30
\bottomrule
\end{tabular}
\caption{Sub-task 1a on the training split. Reranking was evaluated on a seeded random 896 of the
1,211 queries because the GPU session ended first.}
\label{tab:main}
\end{table}

\paragraph{Dense retrieval.} Dense retrieval exceeds BM25 by $+0.1622$ [$+0.1472$, $+0.1774$] on
1a and raises recall@100 from 0.268 to 0.648. The two retrievers' top-100 lists overlap by only
0.060, but almost all of the disagreement goes one way: dense finds 1,706 judged documents that
BM25 misses and BM25 finds 104 that dense misses, so their union adds only 0.027 recall@100.
We therefore did not use BM25.

\subsection{Step fusion is a candidate-pool union}
\label{sec:fusion}
Fusion improves 1a by $+0.0115$ [$+0.0071$, $+0.0162$]. Three measurements locate the gain:

\emph{Weight.} Dropping the whole-query list ($w{=}0$) gives $+0.0014$, not significant. A weight of
0.001 gives the full $+0.0115$. From there the gain is flat to $w{\approx}0.25$ and declines as
$w$ grows ($+0.0030$ at $w{=}5$, not significant at $w{=}10$). A weight too small to reorder
anything that both lists contain still admits the whole query's own candidates into the pool.

\emph{Recall.} Recall@100 is 0.6479 for the whole-query list, 0.6466 for the step lists, and
0.6849 for their union: the two find the same amount of relevant material but different
documents.

\emph{Operator.} Across sum, max and reciprocal rank fusion \citep{cormack2009rrf} with max,
min-max and rank normalisation, every configuration that includes the whole-query list scores
0.2654--0.2695, none significantly different from our choice. The only significant difference
is reciprocal rank fusion at equal weight, $-0.0060$ [$-0.0093$, $-0.0031$], which recovers fully
at $w{=}0.001$. The operator ablation that decomposition work usually reports
\citep{zhong2025redi} matters little here; which lists enter the pool matters.

\emph{Direction.} Adding the whole-query list to each step's list and scoring 1b never helps:
$+0.0000$ at $w{=}0.001$ and a monotone decline to $-0.0083$ [$-0.0120$, $-0.0047$] at $w{=}1$.
A step's own top 10 is already full, so the extra candidates either stay below rank 10 or
displace step-specific ones.

\subsection{Reranking}
On the 896-query subset, reranking adds $+0.0508$ [$+0.0362$, $+0.0664$] over fusion. It improves
12 of 13 domains (bitcoin falls from 0.1660 to 0.1538 on 49 queries); per query it wins 390
times and loses 194. Only 2 of 1,792 reranking windows produced an unparseable answer.

\section{Development check}
\label{sec:dev}
\pending{single dev evaluation of the frozen system: 1a and 1b nDCG@10, compared with train.}

\section{Compute}
All GPU work ran on free Kaggle T4 GPUs (16\,GB, no bf16, no FlashAttention). Corpus embedding:
6.0 hours on two GPUs. Reranking: 29.7 seconds per query on two GPUs with vLLM, 7.6 hours for
896 queries; prompts average 6.5k tokens per window, and the run is bound by prompt processing.
Scoring, fusion and all analyses ran on a CPU laptop.

\section{Limitations}
Reranking was measured on 74\% of the training queries. Its cost scales with the number of test
queries, not the corpus. We did not tune the reranking depth. The first stage uses the smallest
DIVER retriever. Query rewriting (TEMPO's strongest reported lever) was not attempted. All
choices were made on training data from the same 13 domains as the test set
\pending{confirm with organizers that test domains match}.

\section*{Models and data}
\texttt{AQ-MedAI/Diver-Retriever-0.6B} (Apache-2.0); \texttt{liuwenhan/reasonrank-7B}, revision
\texttt{3444046} (MIT); scoring with \texttt{pytrec\_eval} 0.5.10 \citep{vangysel2018pytreceval}.
No external documents were indexed, no gold labels or guidance annotations were used at
inference, and no model was trained or fine-tuned. Code: \pending{public repository URL}.

\bibliography{references}

\end{document}
